% !TeX program = xelatex
\documentclass[UTF8,a4paper,fontset=none]{ctexart}

\usepackage{geometry}
\usepackage{setspace}
\usepackage{fancyhdr}
\usepackage{titlesec}
\usepackage{graphicx}
\usepackage{booktabs}
\usepackage{longtable}
\usepackage{array}
\usepackage{xurl}

% Page layout
\geometry{left=2.5cm,right=2.5cm,top=2cm,bottom=2cm}
\setlength{\parindent}{2em}
\setstretch{1.5}
\setlength{\tabcolsep}{3pt}

% Fonts
\setCJKmainfont[AutoFakeBold=2]{华文宋体}
\setCJKsansfont{冬青黑体简体中文}
\setCJKmonofont{冬青黑体简体中文}
\setCJKfamilyfont{zhsong}{华文宋体}
\setCJKfamilyfont{zhhei}{冬青黑体简体中文}
\newcommand{\songti}{\CJKfamily{zhsong}}
\newcommand{\heiti}{\CJKfamily{zhhei}}
\newcolumntype{L}[1]{>{\raggedright\arraybackslash}p{#1}}
\newcommand{\field}[1]{\path{#1}}
\ctexset{
  section = {format=\heiti\zihao{2}},
  subsection = {format=\heiti\zihao{3}},
  subsubsection = {format=\heiti\zihao{4}}
}

% Footer page number (centered, small-five)
\pagestyle{fancy}
\fancyhf{}
\fancyfoot[C]{\songti\zihao{5}\thepage}

\begin{document}
\songti\zihao{-4}

\begin{titlepage}
\thispagestyle{empty}

\begin{center}
  \vspace*{2cm}
  {\heiti\zihao{3} 采摘工智慧管理系统（Picking-Pass）}\\[1.5em]
  {\heiti\zihao{2} 数据库结构说明文档}
\end{center}

\vspace{5cm}

\noindent
文档信息：

\vspace{2cm}

\noindent
\begin{tabular}{@{}ll}
文档名称：& \texttt{pickpass\_db} 数据库结构说明\\
适用版本：& dev（2026-04-02）\\
数据库类型：& MySQL 8.x\\
对应初始化文件：& \texttt{app/backend/picking\_pass\_db\_final.sql}\\
更新日期：& 2026年04月02日
\end{tabular}
\end{titlepage}

\section{文档目的}
本文档用于说明采摘工智慧管理系统在当前 dev 版本下的数据库结构基线，重点覆盖：核心业务表、字段定义、关键外键关系、约束策略与初始化流程。文档目标是为开发、联调与测试提供统一的数据结构依据，避免“代码已更新但数据库结构未同步”的环境偏差。

\section{数据库设计原则}
数据库设计围绕“用工业务全流程闭环”展开，遵循以下原则：

\noindent 1) 业务主链可追溯：以“用户-岗位-签到-结算-支付”为主链，关键操作写入操作日志。\\
\noindent 2) 敏感信息最小暴露：手机号、身份证、银行卡等字段采用密文存储，配合 Hash 字段支持精确检索。\\
\noindent 3) 约束前置到库层：使用外键、唯一索引、检查约束、触发器保证数据一致性。\\
\noindent 4) 面向联调的可维护性：初始化 SQL、迁移脚本与硬化脚本协同，确保新环境可重复落地。

技术上下文：后端基于 NestJS + TypeORM，数据库为 MySQL。当前文档内容与初始化文件 \texttt{picking\_pass\_db\_final.sql} 对齐。

\subsection{数据库表清单}
根据当前实际数据库 \texttt{pickpass\_db} 的在线结构，系统现有 10 张核心业务表，具体如下：

\begin{tabular}{p{4.2cm}p{9.3cm}}
\toprule
表名 & 主要作用 \\
\midrule
\texttt{sys\_user} & 用户表，存储工人、基地管理员、现场管理员、超级管理员等角色信息，以及手机号、身份证、紧急联系人等敏感字段的密文与 Hash 索引。 \\
\texttt{base\_info} & 基地信息表，存储基地名称、区域、联系方式、营业执照、审核状态、所属负责人等基础信息。 \\
\texttt{recruitment\_job} & 岗位招聘表，存储岗位名称、所属基地、招聘人数、薪资类型、工时要求、福利与有效期等招聘信息。 \\
\texttt{daily\_signup} & 报名与签到表，记录工人按日期对应的岗位报名、签到状态、签到时间、代报名与离线同步情况。 \\
\texttt{labor\_salary} & 工资结算表，记录每次签到对应的工时、计件数、单价快照、总金额、发放状态及经办管理员。 \\
\texttt{salary\_payment} & 工资发放表，记录工资确认、发放方式、支付凭证、确认签字、发放时间与发放人信息。 \\
\texttt{job\_application} & 岗位申请表，记录工人对岗位的申请、审核状态、备注、拒绝原因及审核时间。 \\
\texttt{base\_cooperation} & 基地合作申请表，记录合作申请人、目标基地、合作需求、审核结果与审核信息。 \\
\shortstack[l]{\texttt{offline\_attendance\_}\\\texttt{event}} & 离线考勤事件表，记录离线签到原始事件、风控校验、人工复核与落库结果。 \\
\texttt{operation\_log} & 操作日志表，记录用户在系统中的创建、修改、审核、支付、签到等关键操作，用于审计追溯。 \\
\bottomrule
\end{tabular}

各表之间的核心关系为：\texttt{sys\_user} 与 \texttt{base\_info} 通过负责人字段关联，\texttt{base\_info} 与 \texttt{recruitment\_job} 构成“一基地对多岗位”，\texttt{daily\_signup} 连接用户、基地与岗位，\texttt{labor\_salary} 基于签到记录生成，\texttt{salary\_payment} 再对工资记录进行支付确认；同时 \texttt{offline\_attendance\_event} 通过外键关联用户、基地、岗位与签到记录，用于承载离线签到的补录与审核链路。

\subsection{表结构说明}
下面按照当前数据库中的真实表结构，对各张表的字段进行完整说明。

\subsubsection{\texttt{sys\_user}（用户表）}
\small
\begin{longtable}{@{}L{2.8cm}L{2.4cm}L{3.0cm}L{6.0cm}@{}}
\toprule
字段名 & 数据类型 & 约束 & 含义 \\
\midrule
\endfirsthead
\toprule
字段名 & 数据类型 & 约束 & 含义 \\
\midrule
\endhead
\field{id} & bigint & 主键，自增 & 用户主键 ID。 \\
\field{uid} & varchar(32) & 非空，唯一 & 对外展示的用户唯一编号。 \\
\field{name} & varchar(50) & 非空 & 用户真实姓名。 \\
\field{id_card_enc} & varchar(256) & 非空 & 身份证号密文。 \\
\field{phone_enc} & varchar(256) & 非空 & 手机号密文。 \\
\field{id_card_hash} & varchar(64) & 非空，唯一索引 & 身份证号 Hash，用于检索与去重。 \\
\field{role_key} & enum & 非空，默认 \field{worker} & 用户角色，包含 \field{super\_admin}、\field{region\_admin}、\field{boss}、\field{base\_manager}、\field{field\_manager}、\field{worker}。 \\
\field{face_img_url} & varchar(255) & 可空 & 人脸或头像图片地址。 \\
\field{region_code} & int & 可空 & 所属区域代码。 \\
\field{is_deleted} & tinyint & 非空，默认 0 & 逻辑删除标记。 \\
\field{created_at} & datetime(6) & 非空，默认当前时间 & 创建时间。 \\
\field{updated_at} & datetime(6) & 非空，自动更新 & 更新时间。 \\
\field{phone_hash} & varchar(64) & 可空，唯一索引 & 手机号 Hash，用于登录检索与去重。 \\
\field{assigned_base_id} & bigint & 可空，外键 & 现场管理员对应的基地 ID。 \\
\field{emergency_contact_enc} & varchar(256) & 可空 & 紧急联系人信息密文。 \\
\field{emergency_phone_enc} & varchar(256) & 可空 & 紧急联系人电话密文。 \\
\field{emergency_phone_hash} & varchar(64) & 可空，索引 & 紧急联系人电话 Hash。 \\
\field{info_audit_status} & tinyint & 非空，默认 1 & 信息审核状态：0 待审，1 通过，2 拒绝。 \\
\field{home_address_enc} & varchar(512) & 可空 & 家庭住址密文字段。 \\
\field{bank_name} & varchar(100) & 可空 & 开户银行名称。 \\
\field{bank_card_no_enc} & varchar(256) & 可空 & 银行卡号密文。 \\
\field{bank_card_no_hash} & varchar(64) & 可空，索引 & 银行卡号 Hash。 \\
\bottomrule
\end{longtable}
\normalsize

\subsubsection{\texttt{base\_info}（基地信息表）}
\small
\begin{longtable}{@{}L{2.8cm}L{2.4cm}L{3.0cm}L{6.0cm}@{}}
\toprule
字段名 & 数据类型 & 约束 & 含义 \\
\midrule
\endfirsthead
\toprule
字段名 & 数据类型 & 约束 & 含义 \\
\midrule
\endhead
\field{id} & bigint & 主键，自增 & 基地主键 ID。 \\
\field{base_name} & varchar(100) & 非空 & 基地名称。 \\
\field{license_enc} & varchar(512) & 非空 & 营业执照或资质图片密文。 \\
\field{contact_enc} & varchar(256) & 非空 & 联系电话密文。 \\
\field{category} & tinyint & 非空，默认 1 & 基地类别，如水果、蔬菜等。 \\
\field{region_code} & int & 非空 & 基地所属区域代码。 \\
\field{address} & text & 可空 & 基地详细地址。 \\
\field{description} & text & 可空 & 基地描述信息，可存 JSON 文本。 \\
\field{audit_status} & tinyint & 非空，默认 0 & 审核状态：0 待审，1 通过，2 拒绝。 \\
\field{owner_id} & bigint & 非空，外键 & 基地负责人用户 ID。 \\
\field{is_deleted} & tinyint & 非空，默认 0 & 逻辑删除标记。 \\
\field{created_at} & datetime(6) & 非空，默认当前时间 & 创建时间。 \\
\field{updated_at} & datetime(6) & 非空，自动更新 & 更新时间。 \\
\bottomrule
\end{longtable}
\normalsize

\subsubsection{\texttt{recruitment\_job}（岗位招聘表）}
\small
\begin{longtable}{@{}L{2.8cm}L{2.4cm}L{3.0cm}L{6.0cm}@{}}
\toprule
字段名 & 数据类型 & 约束 & 含义 \\
\midrule
\endfirsthead
\toprule
字段名 & 数据类型 & 约束 & 含义 \\
\midrule
\endhead
\field{id} & bigint & 主键，自增 & 岗位主键 ID。 \\
\field{base_id} & bigint & 非空，外键 & 所属基地 ID。 \\
\field{pay_type} & tinyint & 非空，默认 1 & 薪资类型：固定、时薪、计件。 \\
\field{unit_price} & decimal(10,2) & 可空 & 单价。 \\
\field{targetCount} & int & 可空，默认 0 & 目标数量。 \\
\field{requirements} & text & 可空 & 招聘要求。 \\
\field{status} & tinyint & 非空，默认 1 & 岗位状态。 \\
\field{valid_until} & datetime & 可空 & 有效期截止时间。 \\
\field{created_at} & datetime(6) & 非空，默认当前时间 & 创建时间。 \\
\field{updated_at} & datetime(6) & 非空，自动更新 & 更新时间。 \\
\field{recruit_count} & int & 非空，默认 1 & 招聘人数。 \\
\field{work_cycle} & tinyint & 非空，默认 1 & 结算或工作周期。 \\
\field{work_content} & text & 可空 & 工作内容。 \\
\field{work_hours} & varchar(50) & 可空 & 工作时间段。 \\
\field{work_start_date} & date & 可空 & 开工日期。 \\
\field{work_end_date} & date & 可空 & 截止日期。 \\
\field{salary_amount} & decimal(10,2) & 可空 & 固定工资金额。 \\
\field{hourly_rate} & decimal(10,2) & 可空 & 时薪。 \\
\field{min_age} & tinyint & 可空 & 最小年龄。 \\
\field{max_age} & tinyint & 可空 & 最大年龄。 \\
\field{experience_required} & text & 可空 & 经验要求。 \\
\field{physical_requirement} & text & 可空 & 体力要求。 \\
\field{benefits} & text & 可空 & 福利描述。 \\
\field{has_accommodation} & tinyint & 非空，默认 0 & 是否包住宿。 \\
\field{has_meals} & tinyint & 非空，默认 0 & 是否包餐。 \\
\field{has_transportation} & tinyint & 非空，默认 0 & 是否提供交通补贴。 \\
\field{transportation_subsidy} & decimal(10,2) & 可空 & 交通补贴金额。 \\
\field{workplace_images} & json & 可空 & 工作场景图片列表。 \\
\field{video_url} & varchar(500) & 可空 & 视频地址。 \\
\field{is_active} & tinyint & 非空，默认 1 & 是否有效。 \\
\field{auto_renew} & tinyint & 非空，默认 0 & 是否自动续期。 \\
\field{renewal_days} & int & 非空，默认 7 & 续期天数。 \\
\field{applicant_count} & int & 非空，默认 0 & 已申请人数。 \\
\field{view_count} & int & 非空，默认 0 & 浏览次数。 \\
\field{job_title} & varchar(100) & 非空 & 岗位名称。 \\
\bottomrule
\end{longtable}
\normalsize

\subsubsection{\texttt{daily\_signup}（报名签到表）}
\small
\begin{longtable}{@{}L{2.8cm}L{2.4cm}L{3.0cm}L{6.0cm}@{}}
\toprule
字段名 & 数据类型 & 约束 & 含义 \\
\midrule
\endfirsthead
\toprule
字段名 & 数据类型 & 约束 & 含义 \\
\midrule
\endhead
\field{id} & bigint & 主键，自增 & 报名签到记录 ID。 \\
\field{user_id} & bigint & 非空，外键 & 工人用户 ID。 \\
\field{base_id} & bigint & 非空，外键 & 基地 ID。 \\
\field{job_id} & bigint & 非空，外键 & 岗位 ID。 \\
\field{work_date} & date & 非空 & 工作日期。 \\
\field{status} & tinyint & 非空，默认 0 & 状态：已报名、已签到、缺勤、取消。 \\
\field{checkin_time} & datetime & 可空 & 实际签到时间。 \\
\field{is_proxy} & tinyint & 非空，默认 0 & 是否代报名。 \\
\field{proxy_user_id} & bigint & 可空，外键 & 代报名人用户 ID。 \\
\field{is_offline_sync} & tinyint & 非空，默认 0 & 是否离线同步导入。 \\
\field{created_at} & datetime(6) & 非空，默认当前时间 & 创建时间。 \\
\field{updated_at} & datetime(6) & 非空，自动更新 & 更新时间。 \\
\bottomrule
\end{longtable}
\normalsize

\subsubsection{\texttt{labor\_salary}（工资结算表）}
\small
\begin{longtable}{@{}L{2.8cm}L{2.4cm}L{3.0cm}L{6.0cm}@{}}
\toprule
字段名 & 数据类型 & 约束 & 含义 \\
\midrule
\endfirsthead
\toprule
字段名 & 数据类型 & 约束 & 含义 \\
\midrule
\endhead
\field{id} & bigint & 主键，自增 & 工资记录 ID。 \\
\field{signup_id} & bigint & 非空，唯一，外键 & 对应签到记录 ID。 \\
\field{work_duration} & decimal(4,1) & 非空，默认 0.0 & 工时。 \\
\field{piece_count} & int & 非空，默认 0 & 计件数量。 \\
\field{unit_price_snapshot} & decimal(10,2) & 非空 & 结算时单价快照。 \\
\field{total_amount} & decimal(10,2) & 非空 & 工资总额。 \\
\field{payout_type} & tinyint & 可空 & 发放方式类型。 \\
\field{status} & tinyint & 非空，默认 0 & 结算状态：待确认、已确认、已支付。 \\
\field{proof_img_url} & varchar(255) & 可空 & 证明图片地址。 \\
\field{worker_sign_url} & varchar(255) & 可空 & 工人签字图片地址。 \\
\field{admin_id} & bigint & 非空，外键 & 经办管理员 ID。 \\
\field{created_at} & datetime(6) & 非空，默认当前时间 & 创建时间。 \\
\field{updated_at} & datetime(6) & 非空，自动更新 & 更新时间。 \\
\bottomrule
\end{longtable}
\normalsize

\subsubsection{\texttt{offline\_attendance\_event}（离线考勤事件表）}
\small
\begin{longtable}{@{}L{2.8cm}L{2.4cm}L{3.0cm}L{6.0cm}@{}}
\toprule
字段名 & 数据类型 & 约束 & 含义 \\
\midrule
\endfirsthead
\toprule
字段名 & 数据类型 & 约束 & 含义 \\
\midrule
\endhead
\field{id} & bigint & 主键，自增 & 离线事件主键 ID。 \\
\field{offline_record_id} & varchar(64) & 非空 & 设备侧离线记录 ID（幂等键组成部分）。 \\
\field{device_id} & varchar(128) & 非空 & 采集设备 ID。 \\
\field{worker_uid} & varchar(32) & 非空 & 工人 UID（离线侧上送）。 \\
\field{worker_id} & bigint & 可空，外键 & 解析出的工人 ID。 \\
\field{base_id} & bigint & 非空，外键 & 基地 ID。 \\
\field{job_id} & bigint & 可空，外键 & 岗位 ID。 \\
\field{work_date} & date & 非空 & 逻辑工作日。 \\
\field{occurred_at} & datetime & 非空 & 实际发生时间。 \\
\field{submitted_by} & bigint & 非空，外键 & 提交人 ID。 \\
\field{status} & tinyint & 非空，默认 0 & 审核状态：0 待审、1 自动通过、2 人工通过、3 拒绝。 \\
\field{risk_level} & tinyint & 非空，默认 1 & 风险等级：0 低风险、1 高风险。 \\
\field{validation_message} & text & 可空 & 系统校验说明。 \\
\field{evidence_note} & text & 可空 & 人工复核备注。 \\
\field{evidence_json} & text & 可空 & 证据快照 JSON。 \\
\field{payload_json} & text & 可空 & 原始上报载荷 JSON。 \\
\field{reviewed_by} & bigint & 可空，外键 & 审核人 ID。 \\
\field{reviewed_at} & datetime & 可空 & 审核时间。 \\
\field{applied_signup_id} & bigint & 可空，外键 & 对应落库的签到记录 ID。 \\
\field{created_at} & datetime(6) & 非空，默认当前时间 & 创建时间。 \\
\field{updated_at} & datetime(6) & 非空，自动更新 & 更新时间。 \\
\bottomrule
\end{longtable}
\normalsize

\subsubsection{\texttt{salary\_payment}（工资发放表）}
\small
\begin{longtable}{@{}L{2.8cm}L{2.4cm}L{3.0cm}L{6.0cm}@{}}
\toprule
字段名 & 数据类型 & 约束 & 含义 \\
\midrule
\endfirsthead
\toprule
字段名 & 数据类型 & 约束 & 含义 \\
\midrule
\endhead
\field{id} & bigint & 主键，自增 & 支付记录 ID。 \\
\field{salary_id} & bigint & 非空，唯一，外键 & 对应工资记录 ID。 \\
\field{paymentMethod} & enum & 非空 & 发放方式枚举：\field{cash}、\field{transfer}。 \\
\field{status} & tinyint & 非空，默认 0 & 支付状态：0 待确认，1 已确认，2 已发放，3 已取消。 \\
\field{confirm_signature_url} & text & 可空 & 确认签字图片地址。 \\
\field{payment_voucher_url} & text & 可空 & 支付凭证图片地址。 \\
\field{paid_at} & datetime & 可空 & 发放时间。 \\
\field{paid_by} & bigint & 可空，外键 & 发放人用户 ID。 \\
\field{note} & text & 可空 & 备注。 \\
\field{created_at} & datetime(6) & 非空，默认当前时间 & 创建时间。 \\
\field{updated_at} & datetime(6) & 非空，自动更新 & 更新时间。 \\
\bottomrule
\end{longtable}
\normalsize

\subsubsection{\texttt{job\_application}（岗位申请表）}
\small
\begin{longtable}{@{}L{2.8cm}L{2.4cm}L{3.0cm}L{6.0cm}@{}}
\toprule
字段名 & 数据类型 & 约束 & 含义 \\
\midrule
\endfirsthead
\toprule
字段名 & 数据类型 & 约束 & 含义 \\
\midrule
\endhead
\field{id} & bigint & 主键，自增 & 申请记录 ID。 \\
\field{user_id} & bigint & 非空，外键 & 申请人用户 ID。 \\
\field{job_id} & bigint & 非空，外键 & 目标岗位 ID。 \\
\field{base_id} & bigint & 非空，外键 & 目标基地 ID。 \\
\field{status} & tinyint & 非空，默认 0 & 审核状态。 \\
\field{pending_guard} & tinyint & 计算列，唯一约束参与字段 & 待处理状态幂等保护字段。 \\
\field{note} & text & 可空 & 申请备注。 \\
\field{rejectReason} & text & 可空 & 拒绝原因。 \\
\field{reviewed_by} & bigint & 可空 & 审核人 ID。 \\
\field{reviewed_at} & datetime & 可空 & 审核时间。 \\
\field{created_at} & datetime(6) & 非空，默认当前时间 & 创建时间。 \\
\field{updated_at} & datetime(6) & 非空，自动更新 & 更新时间。 \\
\bottomrule
\end{longtable}
\normalsize

\subsubsection{\texttt{base\_cooperation}（基地合作申请表）}
\small
\begin{longtable}{@{}L{2.8cm}L{2.4cm}L{3.0cm}L{6.0cm}@{}}
\toprule
字段名 & 数据类型 & 约束 & 含义 \\
\midrule
\endfirsthead
\toprule
字段名 & 数据类型 & 约束 & 含义 \\
\midrule
\endhead
\field{id} & bigint & 主键，自增 & 合作申请 ID。 \\
\field{base_id} & bigint & 非空，外键 & 目标基地 ID。 \\
\field{applicant_id} & bigint & 非空，外键 & 申请人用户 ID。 \\
\field{requirement} & text & 非空 & 合作需求描述。 \\
\field{status} & tinyint & 非空，默认 0 & 审核状态。 \\
\field{pending_guard} & tinyint & 计算列，唯一约束参与字段 & 待审核状态幂等保护字段。 \\
\field{rejectReason} & text & 可空 & 拒绝原因。 \\
\field{reviewed_by} & bigint & 可空 & 审核人 ID。 \\
\field{reviewed_at} & datetime & 可空 & 审核时间。 \\
\field{created_at} & datetime(6) & 非空，默认当前时间 & 创建时间。 \\
\field{updated_at} & datetime(6) & 非空，自动更新 & 更新时间。 \\
\bottomrule
\end{longtable}
\normalsize

\subsubsection{\texttt{operation\_log}（操作日志表）}
\small
\begin{longtable}{@{}L{2.8cm}L{2.4cm}L{3.0cm}L{6.0cm}@{}}
\toprule
字段名 & 数据类型 & 约束 & 含义 \\
\midrule
\endfirsthead
\toprule
字段名 & 数据类型 & 约束 & 含义 \\
\midrule
\endhead
\field{id} & bigint & 主键，自增 & 日志记录 ID。 \\
\field{operationType} & enum & 非空 & 操作类型枚举：\field{create}、\field{update}、\field{delete}、\field{audit}、\field{login}、\field{checkin}、\field{payment}。 \\
\field{resourceType} & enum & 非空 & 资源类型枚举：\field{user}、\field{base}、\field{job}、\field{signup}、\field{salary}、\field{offline\_event}。 \\
\field{resource_id} & bigint & 非空 & 被操作资源 ID。 \\
\field{user_id} & bigint & 非空 & 操作人用户 ID。 \\
\field{description} & text & 可空 & 操作描述。 \\
\field{before_data} & text & 可空 & 变更前数据。 \\
\field{after_data} & text & 可空 & 变更后数据。 \\
\field{ip_address} & varchar(45) & 可空 & 操作来源 IP。 \\
\field{user_agent} & text & 可空 & 客户端标识信息。 \\
\field{created_at} & datetime(6) & 非空，默认当前时间 & 操作时间。 \\
\bottomrule
\end{longtable}
\normalsize

\section{关系与约束设计说明}
\subsection{核心关系}
\noindent 1) \texttt{base\_info.owner\_id} 关联 \texttt{sys\_user.id}，限定基地负责人。\\
\noindent 2) \texttt{recruitment\_job.base\_id} 关联基地，形成“一基地多岗位”。\\
\noindent 3) \texttt{daily\_signup} 同时关联用户、基地、岗位，作为考勤与结算桥接表。\\
\noindent 4) \texttt{labor\_salary.signup\_id} 对签到记录一对一结算。\\
\noindent 5) \texttt{salary\_payment.salary\_id} 对工资记录一对一发放。\\
\noindent 6) \texttt{offline\_attendance\_event} 关联用户/基地/岗位/签到，承载离线补录链路。

\subsection{一致性与安全约束}
\noindent 1) 唯一约束：\texttt{sys\_user.uid}、\texttt{id\_card\_hash}、\texttt{phone\_hash} 等用于去重。\\
\noindent 2) 幂等约束：\texttt{job\_application.pending\_guard} 与 \texttt{base\_cooperation.pending\_guard} 防止重复待审。\\
\noindent 3) 触发器约束：限制基地负责人角色、限制 \texttt{assigned\_base\_id} 的写入角色范围。\\
\noindent 4) 敏感字段：\texttt{phone\_enc}、\texttt{id\_card\_enc}、\texttt{bank\_card\_no\_enc} 等采用密文存储，检索使用 Hash 字段。

\section{初始化与同步流程}
\subsection{新环境初始化}
\noindent 1) 创建数据库并导入：\texttt{app/backend/picking\_pass\_db\_final.sql}。\\
\noindent 2) 配置后端 \texttt{.env}（DB 与 AES 参数）。\\
\noindent 3) 启动后端后执行必要硬化脚本（如上线前硬化）。

\subsection{登录基线校验（推荐）}
为避免测试环境出现“手机号可见但登录查不到用户”，建议执行：

\begin{verbatim}
npm run login:baseline
npm run login:baseline:check
\end{verbatim}

\section{维护建议}
\noindent 1) 每次结构变更后，同时更新：初始化 SQL、迁移脚本、本文档。\\
\noindent 2) 联调前先做口径检查（尤其 \texttt{phone\_hash} 与 \texttt{id\_card\_hash}）。\\
\noindent 3) 发布前至少做一次备份与恢复演练，确保可回滚。

\end{document}
